% ch23: 微积分基本定理
\chapter{微积分基本定理——牛顿-莱布尼茨公式}
\begin{applicationbox}
学完本章：掌握积分上限函数、微积分基本定理（N-L公式）
\end{applicationbox}

\section{积分上限函数} \(\Phi(x)=\int_a^x f(t)dt\)，则 \(\Phi'(x)=f(x)\)。

\textbf{ε-δ 证明：} 对任意 \(x\)，取 \(h\neq0\) 使 \(x+h\in(a,b)\)。
\[
\frac{\Phi(x+h)-\Phi(x)}{h} = \frac1h\int_x^{x+h} f(t)dt
\]
由积分中值定理，存在 \(\xi\) 在 \(x\) 与 \(x+h\) 之间使上式 \(=f(\xi)\)。
由于 \(f\) 连续，对任意 \(\varepsilon>0\)，存在 \(\delta>0\) 使 \(|t-x|<\delta\) 时 \(|f(t)-f(x)|<\varepsilon\)。
取 \(|h|<\delta\)，则 \(|\xi-x|<\delta\)，故
\[
\left|\frac{\Phi(x+h)-\Phi(x)}{h} - f(x)\right| = |f(\xi)-f(x)| < \varepsilon
\]
由 ε-δ 定义，\(\Phi'(x)=f(x)\)。\(\blacksquare\)

\section{牛顿-莱布尼茨公式} 若 \(F\) 是 \(f\) 的原函数（\(F'=f\)），则：
\[\int_a^b f(x)dx = F(b)-F(a)\]

\textbf{证明：} 令 \(\Phi(x)=\int_a^x f(t)dt\)，则 \(\Phi'(x)=f(x)=F'(x)\)，
故 \(\Phi(x)-F(x)=C\)（常数）。取 \(x=a\) 得 \(\Phi(a)-F(a)=C\)，即 \(0-F(a)=C\)。
取 \(x=b\) 得 \(\Phi(b)-F(b)=C\)，即 \(\int_a^b f(t)dt - F(b) = -F(a)\)。移项即得。\(\blacksquare\)

\begin{examplebox}[1] \(\int_0^1 x^2 dx = \left[\frac{x^3}{3}\right]_0^1 = \frac13\)\end{examplebox}
\begin{examplebox}[2] \(\int_0^\pi \sin x dx = [-\cos x]_0^\pi = (-\cos\pi)-(-\cos0)=1+1=2\)\end{examplebox}

\section{不定积分} 全体原函数：\(\int f(x)dx = F(x)+C\)

\section{习题}
\begin{exercise}[0] N-L公式的条件是什么？\end{exercise}
\begin{exercise}[0] \(\Phi(x)=\int_0^x t^2 dt\)，求 \(\Phi'(x)\)。\end{exercise}
\begin{exercise}[0] \(\int_a^b f(x)dx\) 用原函数怎么表示？\end{exercise}
\begin{exercise}[0] 不定积分为什么加常数 \(C\)？\end{exercise}
\begin{exercise}[1] 求 \(\int_0^2 x dx\)。\end{exercise}
\begin{exercise}[1] 求 \(\int_0^1 x^3 dx\)。\end{exercise}
\begin{exercise}[1] 求 \(\int_0^{\pi/2} \cos x dx\)。\end{exercise}
\begin{exercise}[1] 求 \(\int_1^e \frac1x dx\)。\end{exercise}
\begin{exercise}[2] 求 \(\int_0^1 (x^2+x) dx\)。\end{exercise}
\begin{exercise}[2] 求 \(\int_{-1}^1 (x^3+2x) dx\)。\end{exercise}
\begin{exercise}[2] 求 \(\int_0^2 |x-1| dx\)。\end{exercise}
\begin{exercise}[2] 求 \(\frac{d}{dx}\int_0^{x^2} \sin t dt\)。\end{exercise}
\begin{exercise}[3] 求 \(\int_0^{\pi} (x+\sin x) dx\)。\end{exercise}
\begin{exercise}[3] 求 \(\int_1^4 \sqrt{x} dx\)。\end{exercise}
\begin{exercise}[3] 求 \(\frac{d}{dx}\int_x^{x^2} e^{t^2} dt\)。\end{exercise}
\begin{exercise}[3] 求 \(\int_0^1 \frac{1}{1+x^2} dx\)。\end{exercise}
\begin{exercise}[3] 求 \(\int_0^{\pi} \sin^2 x dx\)。\end{exercise}
\begin{exercise}[4] 求 \(\int_0^2 f(x)dx\) 其中 \(f(x)=\begin{cases}x,&x\le1\\2-x,&x>1\end{cases}\)。\end{exercise}
\begin{exercise}[4] 证明 \(\frac{d}{dx}\int_a^x f(t)dt = f(x)\)。\end{exercise}
\begin{exercise}[4] 求 \(\int_0^{\pi/2} \frac{\sin x}{1+\cos^2 x} dx\)。\end{exercise}
\begin{exercise}[4] 求 \(\int_0^1 \frac{x}{1+x^2} dx\)。\end{exercise}
\begin{exercise}[4] 已知 \(\int_0^x f(t)dt = x^3\)，求 \(f(x)\)。\end{exercise}
\begin{exercise}[5] 求 \(\int_0^1 \frac{\arctan x}{1+x^2} dx\)。\end{exercise}
\begin{exercise}[5] 证明 \(\int_0^x (\int_0^t f(s)ds)dt = \int_0^x (x-t)f(t)dt\)。\end{exercise}
\begin{exercise}[5] 设 \(F(x)=\int_0^x e^{t^2}dt\)，求 \(F'(x)\) 和 \(F''(x)\)。\end{exercise}
\begin{exercise}[5] 求 \(\lim_{x\to0}\frac{\int_0^x \sin t^2 dt}{x^3}\)。\end{exercise}
\begin{exercise}[6] 证明积分中值定理：\(\exists\xi\in[a,b]\) 使 \(\int_a^b f(x)dx = f(\xi)(b-a)\)。\end{exercise}
\begin{exercise}[6] 求 \(\int_0^1 \frac{\ln(1+x)}{1+x^2} dx\)。\end{exercise}
\begin{exercise}[6] 证明 \(\int_0^x f(t)dt\) 与 \(\int_0^x f(x-t)dt\) 的关系。\end{exercise}
\begin{exercise}[7] 证明：若 \(f\) 连续且 \(\int_a^x f(t)dt=0\) 对所有 \(x\) 成立，则 \(f\equiv0\)。\end{exercise}
